% Cover letter, Physical Review Letters.
% NOT sent, and see the note at the end: we do not currently judge this result
% to meet the PRL criterion. This file exists so the decision is explicit.
\documentclass[11pt]{letter}
\usepackage[margin=1in]{geometry}
\usepackage{amsmath,amssymb}
\newcommand{\TODO}[1]{\textbf{[#1]}}
\signature{\TODO{author names not assigned}}
\address{\TODO{affiliation not assigned}}
\begin{document}
\begin{letter}{The Editors\\ Physical Review Letters}
\opening{Dear Editors,}

\textbf{Internal note, not for submission.} This letter is retained only to
record a judgement made deliberately rather than by omission: we do \emph{not}
consider the present result appropriate for PRL, and the manuscript is directed
to Physical Review D.

The strongest component is a structural no-go: the quasinormal pencil of the
doubly rotating five-dimensional Myers--Perry black hole is a direct sum over
azimuthal sectors, so the equal-spin $U(2)$ multiplet degeneracies cannot be
defective, and the enhanced symmetry closes the very channel it appears to open.
The argument is short and exact, and it rules out a natural expectation.

Against that: the accompanying result is a \emph{bounded negative}
classification over a declared parameter domain, not a discovery; the no-go,
while exact, follows in two lines from a standard property of direct sums once
the decomposition is set up correctly, so its novelty lies in the application
rather than the mathematics; and the one regime where an exceptional point could
still hide in this system --- the quasiresonant limit --- is precisely the
regime our methods cannot reach.

A PRL submission would require either a verified exceptional structure, or a
demonstration that the termination law we derive is realized somewhere, neither
of which we have. Should the hyperboloidal follow-up we identify as the next
step produce either, this assessment should be revisited.

\closing{Sincerely,}
\end{letter}
\end{document}
